% !TeX root = ../navier-stokes-manuscript.tex
\subsection{The local construction behind the maximal history}
\label{app:LocalConstruction}

Local existence and restart use one construction.
Its lifespan depends only on the fixed viscosity and the \(H^3\) size of the datum; once that interval is fixed, every higher Sobolev order already present in the datum persists on it.

For real \(s\), write
\[
  H^s_\sigma(\R^3)
  :=\{v\in H^s(\R^3;\R^3):\nabla\cdot v=0\}.
\]

\begin{proposition}[Local construction and continuation alternative]\label{prop:LocalMaximalHistory}
Let \(t_0\geq0\), \(\nu>0\), and \(a\in H^3_\sigma(\R^3)\).
There is a time \(\delta=\delta(\nu,\norm{a}_{H^3})>0\) and a unique normalized-pressure solution \((w,q)\) with
\[
  w\in C([t_0,t_0+\delta];H^3_\sigma)
  \cap L^2(t_0,t_0+\delta;H^4),
  \qquad
  w_t\in L^2(t_0,t_0+\delta;H^2).
\]
If \(a\in H^m_\sigma\) for an integer \(m>3\), the same interval also carries
\[
  w\in C([t_0,t_0+\delta];H^m_\sigma)
  \cap L^2(t_0,t_0+\delta;H^{m+1}),
  \qquad
  w_t\in L^2(t_0,t_0+\delta;H^{m-1}).
\]
The local solutions join uniquely into one maximal history, and a finite maximal time forces its \(H^3\) norm to become unbounded along the terminal approach.
\end{proposition}

\begin{proof}[Proof of Proposition~\ref{prop:LocalMaximalHistory}]
Fix \(t_0\geq0\), \(\nu>0\), and \(a\in H^3_\sigma(\R^3)\). For real \(s\),
we use the inhomogeneous norm
\[
  \norm{f}_{H^s}^2
  :=\int_{\R^3}\langle\xi\rangle^{2s}
       |\widehat f(\xi)|^2\,d\xi,
  \qquad
  \langle\xi\rangle=(1+|\xi|^2)^{1/2}.
\]
The Leray projection is denoted by \(\Leray\), and \(R_j\) denotes the Riesz
transform with symbol \(i\xi_j/|\xi|\). Both are bounded on every \(H^s\).

\divider{Why the base class has exactly these three rows}

The nonlinear term fixes the base spatial index. In three dimensions,
\[
  H^3(\R^3)\hookrightarrow W^{1,\infty}(\R^3),
  \qquad
  \norm{fg}_{H^3}\leq C\norm{f}_{H^3}\norm{g}_{H^3}.
\]
Thus an \(H^3\) velocity is Lipschitz and \(H^3\) is an algebra. In
particular,
\begin{equation}
  \norm{\Leray\nabla\!\cdot(f\otimes g)}_{H^2}
  \leq C\norm{f}_{H^3}\norm{g}_{H^3}.
  \label{eq:LocalProductEstimate}
\end{equation}
For integer derivatives this follows directly from Leibniz' rule and the
Sobolev embeddings \(H^2\hookrightarrow L^\infty\) and
\(H^1\hookrightarrow L^6\); density gives the stated estimate on \(H^3\).

The heat operator determines the two neighbouring rows. A force in \(H^2\)
produces one square-integrable parabolic derivative above the base state,
\(z\in L^2_tH^4_x\), and the equation places \(z_t\) back in
\(L^2_tH^2_x\). These rows meet exactly at \(H^3\). Indeed, whenever
\[
  v\in L^2(J;H^4),
  \qquad
  v_t\in L^2(J;H^2),
\]
frequency truncation, the fundamental theorem of calculus, and
Cauchy--Schwarz give
\begin{equation}
  \left|
    \norm{v(t)}_{H^3}^2-\norm{v(r)}_{H^3}^2
  \right|
  \leq
  2\norm{v_t}_{L^2(r,t;H^2)}
   \norm{v}_{L^2(r,t;H^4)}.
  \label{eq:LocalAdjacentTrace}
\end{equation}
The right side tends to zero on shrinking time intervals. Passing from the
frequency truncations shows that \(v\) has a unique representative in
\(C(J;H^3)\). This is why the local strong class is
\[
  C_tH^3_x\cap L^2_tH^4_x,
  \qquad
  \partial_t v\in L^2_tH^2_x.
\]
The parabolic gain \(v\in L^2_tH^4_x\) and the equation row
\(v_t\in L^2_tH^2_x\) meet at the continuous \(H^3\) trace needed to carry the
initial value.

\divider{The shifted linear estimate}

Let \(J_\delta=[t_0,t_0+\delta]\), where \(0<\delta\leq1\), and consider
\[
  z_t-\nu\Delta z=F,
  \qquad
  z(t_0)=a,
  \qquad
  F\in L^2(J_\delta;H^2).
\]
The \(H^3\) energy pairs the force with \(z\) through the adjacent spaces
\(H^2\) and \(H^4\). In Fourier variables define
\[
  \langle F,z\rangle_{H^2,H^4;H^3}
  :=\int_{\R^3}\langle\xi\rangle^6
       \widehat F(\xi)\cdot\overline{\widehat z(\xi)}\,d\xi.
\]
The factorization
\(
  \langle\xi\rangle^6
  =\langle\xi\rangle^2\langle\xi\rangle^4
\)
and Cauchy--Schwarz give
\begin{equation}
  \left|\langle F,z\rangle_{H^2,H^4;H^3}\right|
  \leq\norm{F}_{H^2}\norm{z}_{H^4}.
  \label{eq:LocalShiftedPairing}
\end{equation}
For a smooth solution, the \(H^3\) energy identity is
\[
  \frac12\frac d{dt}\norm{z}_{H^3}^2
  +\nu\norm{\nabla z}_{H^3}^2
  =\operatorname{Re}\langle F,z\rangle_{H^2,H^4;H^3}.
\]
The inhomogeneous weights also give the exact splitting
\[
  \norm{z}_{H^4}^2
  =\norm{z}_{H^3}^2+\norm{\nabla z}_{H^3}^2.
\]
Using \eqref{eq:LocalShiftedPairing}, Young's inequality, and then Gronwall's
inequality yields
\[
  \norm{z}_{C(J_\delta;H^3)}
  +\norm{z}_{L^2(J_\delta;H^4)}
  \leq
  C_\nu\left(
    \norm{a}_{H^3}+\norm{F}_{L^2(J_\delta;H^2)}
  \right).
\]
Since \(z_t=\nu\Delta z+F\), the same estimate also controls the time derivative:
\begin{equation}
  \norm{z}_{C(J_\delta;H^3)}
  +\norm{z}_{L^2(J_\delta;H^4)}
  +\norm{z_t}_{L^2(J_\delta;H^2)}
  \leq
  C_\nu\left(
    \norm{a}_{H^3}+\norm{F}_{L^2(J_\delta;H^2)}
  \right).
  \label{eq:LocalLinearEstimate}
\end{equation}
The constant is independent of \(\delta\in(0,1]\). Apply a bounded frequency
cutoff to \(a\) and \(F\), solve the truncated equations, and use
\eqref{eq:LocalLinearEstimate} on differences. The limits solve the general
linear problem, and \eqref{eq:LocalAdjacentTrace} supplies the prescribed
strong \(H^3\) trace at \(t_0\). The same estimate applied to a difference
gives linear uniqueness.

\divider{The projected fixed point}

Set
\[
  X_\delta
  :=C(J_\delta;H^3_\sigma)\cap L^2(J_\delta;H^4),
  \qquad
  \norm{w}_{X_\delta}
  :=\norm{w}_{C H^3}+\norm{w}_{L^2H^4}.
\]
The projected mild equation is
\begin{equation}
  \Phi(w)(t)
  :=e^{\nu(t-t_0)\Delta}a
  -\int_{t_0}^{t}e^{\nu(t-s)\Delta}
    \Leray\nabla\!\cdot(w\otimes w)(s)\,ds.
  \label{eq:LocalMildMap}
\end{equation}
All abbreviated time norms below are taken over \(J_\delta\). Estimate
\eqref{eq:LocalProductEstimate} and the length of the interval give
\begin{align}
  \norm{\Leray\nabla\!\cdot(w\otimes w)}_{L^2H^2}
  &\leq C\delta^{1/2}\norm{w}_{C H^3}^2,
  \label{eq:LocalQuadraticCount}
  \\
  \norm{\Leray\nabla\!\cdot(w\otimes w-v\otimes v)}_{L^2H^2}
  &\leq C\delta^{1/2}
  \bigl(\norm{w}_{C H^3}+\norm{v}_{C H^3}\bigr)
  \norm{w-v}_{C H^3}.
  \label{eq:LocalDifferenceCount}
\end{align}
Inserting these bounds into \eqref{eq:LocalLinearEstimate} gives
\[
  \norm{\Phi(w)}_{X_\delta}
  \leq C_\nu\norm{a}_{H^3}
  +C_\nu C\delta^{1/2}\norm{w}_{X_\delta}^2
\]
and
\[
  \norm{\Phi(w)-\Phi(v)}_{X_\delta}
  \leq C_\nu C\delta^{1/2}
  \bigl(\norm{w}_{X_\delta}+\norm{v}_{X_\delta}\bigr)
  \norm{w-v}_{X_\delta}.
\]
For \(a\neq0\), take \(R=2C_\nu\norm{a}_{H^3}\) and choose
\(\delta\leq1\) so that \(4C_\nu C\delta^{1/2}R\leq1\). The map \(\Phi\)
then preserves the closed \(X_\delta\)-ball of radius \(R\) and has Lipschitz
constant at most \(1/2\) there. If \(a=0\), the fixed point is the zero
solution. Banach's fixed-point theorem produces \(w\in X_\delta\), and the
choice may be made so that
\begin{equation}
  \delta
  \geq
  \min\left\{1,c_\nu(1+\norm{a}_{H^3})^{-2}\right\}.
  \label{eq:LocalLifespan}
\end{equation}
Here \(c_\nu>0\) depends only on the fixed viscosity and the normalization of
the Sobolev estimates.

The contraction first proves uniqueness inside the chosen ball. It also gives
uniqueness in the full stated strong class. If \(w,v\in X_\delta\) have the
same datum, apply \eqref{eq:LocalLinearEstimate} and
\eqref{eq:LocalDifferenceCount} to \(w-v\) on a sufficiently short initial
subinterval. The coefficient is then smaller than one, so \(w=v\) there.
Their common endpoint is the initial value for the next subinterval. A finite
partition of \(J_\delta\) repeats the argument across the entire interval.
The heat semigroup and \(\Leray\) preserve divergence-free fields, and the
projected equation together with \eqref{eq:LocalQuadraticCount} gives
\begin{equation}
  w_t\in L^2(J_\delta;H^2).
  \label{eq:LocalTimeRow}
\end{equation}

\divider{Recovering the normalized pressure}

Define
\begin{equation}
  q:=R_iR_j(w_iw_j).
  \label{eq:LocalPressureFormula}
\end{equation}
Because \(H^3\) is an algebra and the Riesz transforms are bounded on \(H^3\),
\[
  q\in C(J_\delta;H^3).
\]
For \(F=\nabla\!\cdot(w\otimes w)\), the Fourier symbols give
\(\Leray F=F+\nabla q\). The projected fixed point consequently solves
\[
  w_t+(w\cdot\nabla)w+\nabla q=\nu\Delta w,
  \qquad
  \nabla\cdot w=0,
  \qquad
  w(t_0)=a.
\]
The condition \(q(t)\in H^3(\R^3)\) fixes the whole-space pressure constant.
If two such pressures accompany the same velocity, their difference has zero
gradient and belongs to \(L^2\), hence vanishes. The projected uniqueness just
proved is thus uniqueness of the pressure-complete pair \((w,q)\).

\divider{Higher spatial orders on the base interval}

Fix an integer \(m\geq3\) and first suppose \(a\) is smooth. Write
the \(H^m\) energy as the sum of the \(L^2\) energies of
\(\partial^\alpha w\), \(|\alpha|\leq m\). Apply \(\partial^\alpha\) to the
projected equation and pair with \(\partial^\alpha w\). The divergence-free
transport of the top derivative cancels:
\[
  \bigl((w\cdot\nabla)\partial^\alpha w,\partial^\alpha w\bigr)_2=0.
\]
The remaining terms in the Leibniz expansion contain at least one derivative
on the transporting factor. Summing over \(|\alpha|\leq m\), the Sobolev
embeddings \(H^2\hookrightarrow L^\infty\) and \(H^1\hookrightarrow L^6\)
give the integer commutator bound
\[
  \left(
  \sum_{|\alpha|\leq m}
  \norm{[\partial^\alpha,w\cdot\nabla]w}_2^2
  \right)^{1/2}
  \leq C_m\norm{\nabla w}_\infty\norm{w}_{H^m}.
\]
The differentiated energy identity gives
\begin{equation}
  \frac12\frac d{dt}\norm{w}_{H^m}^2
  +\nu\norm{\nabla w}_{H^m}^2
  \leq C_m\norm{\nabla w}_\infty\norm{w}_{H^m}^2.
  \label{eq:LocalPersistenceEnergy}
\end{equation}
The base solution lies in \(C(J_\delta;H^3)\), so
\(\norm{\nabla w}_\infty\) is bounded, and in particular integrable, on this
same interval. Gronwall gives
\begin{equation}
  \norm{w(t)}_{H^m}^2
  \leq
  \norm{a}_{H^m}^2
  \exp\left(
    2C_m\int_{t_0}^{t}\norm{\nabla w(s)}_\infty\,ds
  \right).
  \label{eq:LocalPersistenceBound}
\end{equation}
Keeping the dissipative term in \eqref{eq:LocalPersistenceEnergy} gives
\(w\in L^2(J_\delta;H^{m+1})\). The product estimate at order \(m-1\)
puts the projected transport term in \(L^2(J_\delta;H^{m-1})\), and the
equation gives
\[
  w_t\in L^2(J_\delta;H^{m-1}).
\]
The order-\(m\) version of \eqref{eq:LocalAdjacentTrace} now yields
\begin{equation}
  w\in C(J_\delta;H^m)
  \cap L^2(J_\delta;H^{m+1}),
  \qquad
  w_t\in L^2(J_\delta;H^{m-1}).
  \label{eq:LocalHigherPersistence}
\end{equation}
Formula \eqref{eq:LocalPressureFormula} and the \(H^m\) algebra estimate give
\(q\in C(J_\delta;H^m)\).

For a general \(a\in H^m_\sigma\), approximate \(a\) by divergence-free
frequency cutoffs \(a_k\). Their \(H^3\) norms are bounded by a common
multiple of \(\norm{a}_{H^3}\), so \eqref{eq:LocalLifespan} places every
smooth solution on the same base interval. The fixed-point estimate identifies
their limit with the \(H^3\) solution already constructed, while
\eqref{eq:LocalPersistenceEnergy}--\eqref{eq:LocalPersistenceBound} give
uniform bounds in the three spaces of \eqref{eq:LocalHigherPersistence}.
Weak compactness identifies the higher-order limits with \(w\). The adjacent
trace gives a continuous \(H^m\) representative, and the energy inequality at
\(t_0\), combined with weak convergence to \(a\), gives strong convergence
\(w(t)\to a\) in \(H^m\) as \(t\downarrow t_0\).

No higher order changes the interval: its length was chosen once from \(\nu\)
and \(\norm{a}_{H^3}\). If \(a\in H^\infty_\sigma\), applying
\eqref{eq:LocalHigherPersistence} at each fixed \(m\) and using \(H^3\)
uniqueness shows that all orders belong to the same pair. Repeatedly
differentiating the pressure formula and the momentum equation then gives
smoothness in time with values in every finite Sobolev space. Spatial Sobolev
embedding makes \((w,q)\) smooth in \((x,t)\) on the construction interval.

\divider{The maximal history and its continuation alternative}

Starting at \(t_0\), join all local intervals furnished by the construction.
Strong uniqueness identifies the solutions and normalized pressures on every
overlap, so their union is one pressure-complete solution on a maximal interval
\([t_0,T_{\max})\).

Suppose \(T_{\max}<\infty\) and
\[
  \sup_{t_0\leq t<T_{\max}}\norm{w(t)}_{H^3}\leq M.
\]
The lifespan calculation supplies the same positive lower bound
\[
  \delta_M
  :=\min\left\{1,c_\nu(1+M)^{-2}\right\}
\]
whenever the construction is started from the actual datum \(w(t)\) at a later
time. Choose \(t>T_{\max}-\delta_M/2\). The solution starting from \(w(t)\)
exists beyond \(T_{\max}\), and uniqueness identifies it with the maximal
history on their common interval. This extends that history past its maximal
time, which is impossible. We have proved the continuation alternative
\begin{equation}
  T_{\max}<\infty
  \quad\Longrightarrow\quad
  \limsup_{t\uparrow T_{\max}}\norm{w(t)}_{H^3}=\infty.
  \label{eq:LocalContinuationAlternative}
\end{equation}
Every higher order present in the initial datum persists on each local piece,
because each piece uses the same \(H^3\)-controlled lifespan. For
\(a\in\Ssigma\subset H^\infty_\sigma\), this gives the smooth
pressure-normalized maximal history asserted in the proposition.
\end{proof}
